\heading{数据结构与算法A-16秋-期中}
\noteinfo{题目来源为真题扫描件转录。客观题答案和部分参考解答已整理在解答中；原扫描中考场纪律、装订线等非题目信息已尽量删去。本次整理提供 C 语言版与 Python 语言版；除程序代码语言外，两版题目内容保持一致。}

\begin{question}
一、选择题（每题 2 分，共 20 分）\\
1. 下面函数的时间复杂度是\_\_\_\_\_\_\_\_\_\_\_。\\
int work(int n, int m) \{\\
int sum = 0;\\
for (int i = 1; i <= m; i *= 2)\\
for (int j = 1; j <= n; j *= 2) \{\\
sum += i * j;\\
\}\\
return sum;\\
\}\\
2. 下面关于线性表的叙述中，正确的是哪些？\\
A. 采用顺序存储的线性表，必须占用一片连续的存储单元。\\
B. 采用顺序存储的线性表，便于进行插入和删除操作。\\
C. 采用链接存储的线性表，不必占用一片连续的存储单元。\\
D. 采用链接存储的线性表，便于插入和删除操作。\\
E. 采用链接存储的线性表，插入第i 个元素的时间与i 的数值成正比。\\
F. 采用链接存储的线性表，在已知的结点p 后插入一个新节点的时间复杂度与结点p 的\\
位置有关。\\
3. 若元素a、b、c、d 依次进栈，则可能的出栈序列有\_\_\_\_\_\_\_\_种；若有元素a、b、c、d、\\
e 依次进栈，则可能的出栈序列有\_\_\_\_\_\_\_\_种。\\
4. 一个字符串中\_\_\_\_\_\_\_\_\_\_\_\_\_\_\_\_\_\_\_\_\_称为该串的子串。abcab 的不相等真子串个数\\
为\_\_\_\_\_\_\_\_\_\_\_\_。\\
5. 一个具有n 个结点的二叉树的叶结点的个数最多为\_\_\_\_\_\_\_\_\_\_\_\_\_\_\_\_\_\_\_\_\_\_\_\_。\\
6. 一个有5 层结点的完全3 叉树。按前序遍历周游给结点从1 开始编号，则第100 号结点的\\
父结点的编码为 \_\_\_\_\_\_。（注释：独根树为1 层的）\\
7. 将序列 \{12, 70, 33, 65, 24, 56, 48, 92, 86, 33\} 按建堆方法调整为小值堆，\\
调整后的序列为\_\_\_\_\_\_\_\_\_\_\_\_\_\_\_\_\_\_\_\_\_\_\_\_\_\_\_\_\_\_\_\_。\\
8. 假设用于通信的电文仅由8 个字符 \{a,b,c,d,e,f,g,h\} 组成，字符在电文中出现的频\\
率分别为0.07，0.19，0.02，0.06，0.32，0.03，0.21，0.10，对这些字符进行哈\\
夫曼编码的外部路径长度为\_\_\_\_\_\_\_\_\_\_。\\
9. 将森林F 转换为对应的二叉树T，F 中的叶结点的个数为\_\_\_\_\_\_\_\_\_\_\_\_\_\_\_\_\_\_。\\
A. T 中叶结点的个数\\
B. T 中度为1 的结点个数\\
C. T 中左子指针为空的结点个数\\
D. T 中右子指针为空的结点个数\\
10. 一棵树T 有20 个度为4 的结点，10 个度为3 的结点，1 个度为2 的结点，10 个度为1\\
的结点，则树T 的结点个数为\_\_\_\_\_\_\_\_\_\_\_\_。\\
\\
二、 辨析与简答(共32 分)\\
1. （6 分）请谈谈循环队列和链表的优缺点。当遇到如下情况时，应当使用哪一种数据结构？\\
情况：你要维护食堂鸡腿饭的排队队伍，队伍有两种操作：1. 排在队首的人出来打鸡腿\\
饭后离开； 2. 一个新来的人排在队尾。你知道食堂不太大，所以队伍的最大长度你可以\\
估计出来。\\
2. （4 分）如何找出具有相同的中序序列和后序序列的所有二叉树？\\
3. （4 分）从空二叉树开始，将关键码\{18,73,10,5,68,99,27,41,51,32,25\}严格按\\
照BST（二叉搜索树）的插入算法逐个插入BST， 画出这样构造出的BST 树，并写出对\\
该BST 按照前序遍历得到的序列。\\
4.\\
（6分）以下是计算模式P的next向量的算法：\\
int *findNext(P) \{\\
int i = 0,  k = -1;\\
int m = len(P);\\
int* next = int[m];\\
next[0] = -1\\
while i < m:\\
while (k >= 0 and P[i] != P[k])\\
k = next[k];\\
i++, k++;\\
if i == m: break\\
if (P[i] == P[k])\\
next[i] = next[k];\\
else\\
next[i] = k;\\
\}\\
return next;\\
\}\\
采用上述算法，计算模式P="abcdaabcab"的next向量。\\
i\\
P[i]\\
a\\
b\\
c\\
d\\
a\\
a\\
b\\
c\\
a\\
b\\
next[i]\\
5.\\
（6分）数组\{23, 17, 14, 6 ,13, 10, 1, 5, 8, 12\}是否为一个最大值堆？若是，\\
请说明理由，否则请严格按照筛选法建堆的过程将其调整成为最大值堆，并画出逐步调整\\
建堆的过程。\\
6.\\
（6分）对下列15个等价对进行合并，给出所得等价类树的图示。在初始情况下，集合中\\
的每个元素分别在独立的等价类中。使用重量权衡合并规则，合并时子树结点少的并入结\\
点数多的那棵（多的那个作为新树根，少的那个根作为新根的直接子结点）；若两棵树规\\
模同样大，则把根值较大的并入根值较小（新树根取值小的）。同时，采用路径压缩优化。\\
(0,2) (1,2) (3,4) (3,1) (3,5) (9,11) (12,14) (3,9) (4,14) (6,7) (8,10)\\
(8,7) (7,0) (10,15) (10,13)\\
\\
三、 算法填空(每空2 分，共16 分)\\
1. 下面的算法将一个用带度数的后根次序法表示的森林转换为左子结点/右兄弟结点法表示。\\
请利用题目给出的树结点ADT和栈ADT，填充空格，使其成为完整的算法。空格中可能需要\\
填写0到多条语句（或表达式）。\\
template<class value\_type>\\
class stack \{\\
public:\\
empty();                        // 判断栈空\\
int size();                          // 返回栈大小\\
value\_type \&top();                  // 读栈顶\\
def push(const value\_type\& X);  // 入栈\\
def pop();                          // 出栈\\
\};\\
class Node<T> \{\\
T info;                               // 结点的数据信息\\
int degree;                          // 结点的度数信息\\
\};\\
template<class T>\\
class TreeNode \{                         // 树结点类\\
public:\\
isLeaf();                       // 判断当前结点是否为叶结点\\
T Value();                            // 返回结点的值\\
TreeNode<T> *LeftMostChild();     // 返回第一个子结点\\
TreeNode<T> *RightSibling();      // 返回右兄\\
def setValue(const T\&);           // 设置当前结点的值\\
def setChild(TreeNode<T> *pointer);    // 设置左子\\
def setSibling(TreeNode<T> *pointer);  // 设置右兄\\
\};\\
template<class T>\\
TreeNode<T> *Convert(Node<T>* nodes, int size) \{\\
TreeNode<T> *cur, *temp1, *temp2;\\
stack<TreeNode<T>*> Cstack;\\
for ( int i = 0; i < size; i++) \{\\
cur = TreeNode<T>(nodes[i].info);\\
if ( nodes[i].degree == 0 )\\
// 填空1\\
else \{\\
assert( nodes[i].degree <= Cstack.size() );\\
temp2 = None;\\
for (int j = 0;       // 填空2      ; j++) \{\\
temp1 = Cstack.top();\\
Cstack.pop();\\
// 填空3\\
temp2 = temp1;\\
\}\\
// 填空4\\
Cstack.push(cur);\\
\}\\
\}\\
cur = temp2 = None;\\
while( not Cstack.empty() ) \{\\
cur = Cstack.top();\\
Cstack.pop();\\
// 填空5\\
temp2 = cur;\\
\}\\
return cur;\\
\}\\
2. 给定一个二叉树的根结点和树中任意两个结点，补全下面的程序段以求这两个结点的最近\\
公共祖先。\\
TreeNode* lowestCommonAncestor(\\
TreeNode* root, TreeNode* p, TreeNode* q) \{\\
if (      // 填空6      ) return root;\\
TreeNode* left = lowestCommonAncestor(root.left, p, q);\\
TreeNode* right = lowestCommonAncestor(root.right, p, q);\\
if (\_      // 填空7      \_\_) return root;\\
return \_      // 填空8      \_\_? left:right;\\
\}\\
\\
四、 算法设计与实现 (16 分)\\
1.\\
（8 分） 现有两个满足先进先出原则的队列A 和B，在队列上可进行以下4 个操作：\\
front(): 获得队首元素\\
push\_back(T x)：把元素x 插到队尾。\\
pop\_front(): 删除队首元素。\\
empty(): 判断是否为空。\\
请用A 和B 模拟一个栈的4 个操作：top()，pop()，push()，empty()。只要给出实\\
现栈操作的基本思想，对时间效率无具体要求。\\
2.\\
（8 分） 对于两个串A 和B ，给出一种算法使得能够保证正确地求出最大的L 值，使得\\
A 串中长为L 的前缀与B 串中长为L 的后缀相等。\\
\\
五、 分析证明题 (共 16 分)\\
1.\\
（6 分）试证明在一棵树中，u 是v 的祖先，当且仅当在先序遍历中u 在v 之前且在后序\\
遍历中u 在v 之后。\\
2.\\
（10 分）对于满二叉树：\\
（1） 试问具有n 个叶结点的树中共有多少个结点？\\
（2） 试证明：\\
，其中n 为叶结点的个数，\\
il 表示第i 个叶结点所在的\\
层次（设根结点所在的层次为0）。\\
\\
\

\begin{solution}
本题为整卷转录型资料：选择题、填空题的参考答案以及简答、算法设计题的参考做法已统一放在本解答中，题面不保留原扫描夹带的答案标注。对扫描不清处，保留了原转录中可辨认的信息，后续可继续按原 PDF 图示补充图片。

\medskip
\noindent 原转录中夹带在题面里的参考答案/解答摘录如下，现已移入解答：
本卷包含选择填空、简答、算法设计等题型。下面保留按扫描件整理的题面；其中客观题参考答案已填入空格或括号，主观题给出参考思路或代码。

得分\\

\end{solution}
\end{question}
